\documentclass[12pt]{article}
\usepackage[margin=1in]{geometry}
\usepackage{setspace}
\usepackage{mathptmx}

\begin{document}
\thispagestyle{empty}

\begin{center}
{\large\textbf{$90{+}7'$: Goal Displacement and Norm Coordination\\
After FIFA's Stoppage-Time Directive}}

\bigskip

{[Author Name]}\\
\textit{[Affiliation]}\\
{[email]}

\bigskip

\textit{JEL codes: L83, J44}
\end{center}

\bigskip

\noindent
In November 2022, FIFA instructed referees to fully compensate for time
lost during professional football matches. Using data from 15~European
leagues and 36,000 matches, we document three findings. First, enforcement intensity tracked institutional proximity to
FIFA: the Premier League added 2.1~minutes per match while Serie~A
added only 0.7~minutes despite identical formal treatment, and
within-country comparisons reproduce the gradient. Second, total goals per match were unchanged, but goals
\textit{displaced} from the final minutes of regulation into expanded
stoppage time---the share of result-changing stoppage-time goals rose
from 7.3\% to 10.3\%. A model of optimal attacking intensity under
fatigue constraints rationalizes this: extending the horizon
redistributes effort without increasing output. Third, 95\% of referees
individually increased stoppage time with no decay over three seasons,
consistent with focal-point coordination rather than enforcement.

\end{document}
